\chapter{Formulation：检索式可编辑参数控制}
\label{chap:formulation}

本文将一个效果预设表示为参数向量 $\theta$。在代码中，$\theta$ 表现为嵌套 JSON，而非单纯的稠密向量；叶子节点包含连续参数、字符串形式数值和模块开关状态。评价时，\texttt{Evaluator} 会递归展开数值叶子，忽略无法转为数值的字符串，并在缺失数值键时使用零补齐。因此，本文中的参数距离表示当前实验代码定义的 flattened numeric parameter RMSE；它不等同于感知距离或插件内部物理距离。

令 $q=(x,a)$ 表示一次查询，其中 $x$ 是文本描述，$a$ 是可选参考音频或其缓存向量。令 $\mathcal{D}=\{(x_i,a_i,\theta_i)\}_{i=1}^{N}$ 表示预设知识库。检索方法学习或定义一个相似度函数 $s(q,i)$，并返回
\[
	\hat{i}=\arg\max_i s(q,i), \qquad \hat{\theta}=\theta_{\hat{i}}.
\]
实验评价的核心对象是 $\hat{\theta}$ 与 held-out query 的 ground-truth 参数 $\theta^\star$ 之间的差异。该形式化强调两个事实：本文研究的是 retrieval prior，范围不覆盖从零生成参数；评价目标是参数对齐，范围不覆盖完整听感偏好。

NeurIPS 主文中的问题定义应进一步固定为检索式可编辑参数控制。给定训练侧或知识库侧预设集合 $\mathcal{D}_{KB}$、测试侧查询集合 $\mathcal{D}_{test}$、查询输入 $q=(x,a)$ 和目标参数 $\theta^\star$，方法 $r_\phi$ 直接输出可被插件或效果链执行的候选参数对象，而非新生成波形：
\[
	r_\phi(q,\mathcal{D}_{KB}) \rightarrow (\hat{i},\hat{\theta},\hat{m}), \qquad \hat{\theta}=\theta_{\hat{i}}.
\]
其中 $\hat{m}$ 表示候选预设的模块结构或元数据。主风险函数定义为 held-out 查询上的期望参数误差
\[
	\mathcal{R}(r_\phi)=\mathbb{E}_{(q,\theta^\star)\sim\mathcal{D}_{test}}\left[d(\hat{\theta},\theta^\star)\right],
\]
其中 $d(\cdot,\cdot)$ 在当前代码中由 flattened numeric parameter RMSE、Acc@0.1、Recall、Cosine 和 Module consistency 共同操作化。该 formulation 的关键约束是 executable：候选 $\hat{\theta}$ 必须来自真实预设或可被真实效果链加载的参数 JSON。这个约束区分了本文任务与纯文本到音频生成、纯 embedding retrieval、不可执行的风格标签分类。

因此，主论文的中心假设可以写成以下可证伪命题：在相同知识库、相同 held-out split 和相同参数距离定义下，二阶纹理统计表示比一阶音频 embedding、文本/关键词匹配和通用音频语言 embedding 更适合检索可执行效果参数。若后续 PaSST、PANNs、direct regression 或 Text2FX-style optimization 在同一 logger schema 下超过 TRR，该命题必须被削弱或重写。

\begin{table}[H]
	\centering
	\caption{检索式可编辑参数控制的 formal contract。}
	\label{tab:formulation-contract}
	\small
	\resizebox{\textwidth}{!}{%
	\begin{tabular}{lll}
		\toprule
		对象 & 符号 & 本文约束 \\
		\midrule
		查询 & $q=(x,a)$ & 文本描述与可选参考音频，均可退化或缺失。 \\
		知识库 & $\mathcal{D}_{KB}$ & 只包含可执行预设与可审计参数 JSON。 \\
		检索器 & $r_\phi$ & 输出候选索引、候选参数和候选模块元数据。 \\
		目标 & $\theta^\star$ & held-out query 的 ground-truth 参数。 \\
		风险 & $\mathcal{R}(r_\phi)$ & 在测试查询上最小化参数空间距离。 \\
		证据 & run artifact & 每个主张必须绑定 split、checksum、per-query CSV 和统计检验。 \\
		\bottomrule
	\end{tabular}}
\end{table}

表~\ref{tab:formulation-contract} 将本节符号和后续实验台账连接起来。后续每个实验若改变输入、知识库、split 或风险定义，都必须在 logger 中形成新的 \texttt{run\_id}，不能与既有主表混写。

TRR 指 Texture Resonance Retrieval。它在实验实现中优先读取 \texttt{item["Vectors"]["TRR"]} 或 \texttt{<AudioPath>.trr.npy} 缓存，只有缺失缓存时才懒加载 \texttt{TextureEncoder(project\_dim=64)} 计算嵌入。检索时，TRR 对查询和数据库嵌入做 L2 归一化，并以余弦相似度排序。该实现路径说明，当前实验主要验证缓存 TRR 表示在既定数据集上的检索效果，尚未验证在线特征提取系统的工程吞吐。

文本基线在报告中应谨慎命名。产品侧 RAG 使用 ChromaDB、SentenceTransformer、文本集合和参数集合；但实验侧的 \texttt{RAGRetriever} 会将 SentenceTransformer monkey-patch 为随机 dummy embedding，并在文本检索中实际使用 deterministic keyword overlap。因此，实验中的 Text-RAG 更准确地说是轻量文本/关键词检索基线，尚未构成完整语义向量 RAG 的强实现。

本文使用五个主要客观指标。Param. Dist. 是展开数值叶子后的 RMSE，越低越好。Acc@0.1 是数值叶子处于绝对误差 0.1 内的比例，越高越好。Recall 衡量激活参数的召回行为。Cosine 是展开参数向量的方向相似度。Module 是根据键名中 On/Off 模块状态推导的模块一致性。由于这些指标均来自当前代码的 operational definition，跨协议比较绝对数值时必须谨慎；同一协议内的配对比较更可靠。

最后，本文使用 evidence boundary 指一个结论与其证据来源之间的匹配关系。如果某个结果有 CSV、JSON、统计报告和代码路径共同支撑，它可以进入主证据链。如果某个结果只有报告文字而缺少原始输出，或者来自 deprecated 文件，它只能作为历史线索。这个规则贯穿本文所有结论。
